\newcommand{\nomeEstado}{BAHIA}
\newcommand{\siglaEstado}{BA}
\newcommand{\nomeConcessionaria}{NEOENERGIA COELBA}
\newcommand{\cnpjConcessionaria}{15.139.629/0001-94}
\newcommand{\presidenteConcessionaria}{Sr. Thiago Guth.}

\newcommand{\empresaResponsavel}{ABEL}
\newcommand{\nomeMunicipio}{Dário Meira}
\newcommand{\nomeMunicipioMaisculo}{\MakeUppercase{\nomeMunicipio}}
\newcommand{\cnpjMunicipio}{13.700.174/0001-09}
\newcommand{\sedePrefeitura}{Rua Isaias Rego, nº 01, Centro}
\newcommand{\generoPrefeitura}{pela Prefeita}
\newcommand{\nomePrefeito}{Marivane Dias Santos}
\newcommand{\numeroContrato}{060601/2024} % Número do contrato (Não precisa por "Nº")
\newcommand{\quantidadeAditivos}{3} % Quantidade de aditivos do contrato
\newcommand{\legitimidade}{2} % 0 - Sem legitimidade, C - No contrato, 1,2,3... - No Aditivo (colocar o número do aditivo)
\newcommand{\publicacao}{Diário Oficial do Município}
\newcommand{\dataPublicacao}{8 de Julho de 2024} % Data da publicação
\newcommand{\tituloClausula}{CLÁUSULA PRIMEIRA - DO OBJETO E REGIME DE EXECUÇÃO} % Título da cláusula que confere poderes de representação
\newcommand{\clausula}{CLÁUSULA PRIMEIRA - DO OBJETO E REGIME DE EXECUÇÃO 

O objeto do presente instrumento é a contratação de empresa para prestação de serviços de assessoria e consultoria técnica, jurídico-administrativo especializada na gestão, elaboração de auditorias e laudos técnicos, mediante a conferência das faturas de energia elétrica da administração direta e indireta do Município, elaboração de memorial de cálculo de consumo e potência do parque de iluminação pública, a apuração do modelo tarifário aplicado em cada unidade consumidora, assim como verificação de possíveis isenções indevidas e/ou não repasse da Contribuição de Iluminação Pública (CIP) e/ou não recolhimento do ISS dos prestadores de serviços do setor elétrico, visando a repetição de indébitos decorrentes de cobranças indevidas (a maior) nas contas de energia elétrica de titularidade, “Deverá a empresa efetuar levantamento dos créditos a que faz jus o município, referentes aos pagamentos indevidos à concessionária de energia elétrica, conforme os prazos estabelecidos na Resolução Normativa da ANEEL nº 1.000, de 7 de dezembro de 2021, especialmente o art. 323, §2º, inciso I, e suas alterações posteriores, podendo para tanto, representar administrativamente o Município de Dário Meira - BA perante a Agência Nacional de Energia Elétrica – ANEEL, em todos os processos, reclamações, recursos e demais atos decorrentes de suas competências.} % Texto da cláusula que confere poderes de representação

